\begin{obeactivity}
\textbf{Aktivitas 14.1: Invers Modulo dan Pencarian Akar Kuadrat}
1. Jalankan fungsi pada Listing~\ref{lst:ecc-invers} untuk menghitung
   $3^{-1} \bmod 11$, $8^{-1} \bmod 11$, dan $20^{-1} \bmod 23$.
2. Bandingkan hasilnya dengan perhitungan manual pada Contoh~\ref{ex:gf11}.
3. Gunakan fungsi pencari akar untuk memecahkan $x^2 \equiv 5 \pmod{11}$
   (Persamaan~\eqref{eq:akar-gf11}). Berapa banyak akar yang ditemukan?
4. Sekarang coba pecahkan $x^2 \equiv 3 \pmod{11}$. Apa yang dikembalikan
   fungsi itu, dan apa artinya bagi penyelesaian persamaan kurva?
\end{obeactivity}

Seluruh operasi pada medan berhingga bertumpu pada dua perkakas pada
Listing~\ref{lst:ecc-invers}: invers modulo dengan algoritma Euclidean yang
diperluas, dan pencarian akar kuadrat dengan penelusuran langsung.

\lstinputlisting[
    language=Python,
    caption={Invers modulo dan pencarian akar kuadrat di dalam $GF(p)$},
    label={lst:ecc-invers},
    firstline=27,
    lastline=44
]{\codefile{python/ecc_uji.py}}

\begin{obeactivity}
\textbf{Aktivitas 14.2: Aritmetika Polinomial di dalam $GF(2^m)$}
1. Jalankan fungsi penjumlahan pada Listing~\ref{lst:ecc-gf2m} untuk
   $\code{0D} + \code{06}$ dan $\code{57} + \code{83}$.
2. Jalankan fungsi perkalian untuk $\code{57} \cdot \code{83}$ dengan polinom
   reduksi $\code{11B}$. Periksa bahwa hasilnya $\code{C1}$, seperti pada
   Contoh~\ref{ex:gf2m}.
3. Ulangi untuk $\code{0D} \cdot \code{06}$ dan bandingkan hasilnya dengan
   penjumlahan kedua byte itu.
4. Jalankan fungsi pembagian panjang untuk membagi hasil kali mentah
   $x^{13} + x^{11} + x^9 + x^8 + x^6 + x^5 + x^4 + x^3 + 1$ dengan
   $\code{11B}$. Berapa hasil bagi dan sisanya?
5. Gunakan fungsi pemeriksa ketereduksian untuk menguji $x^4 + x + 1$ dan
   $x^2 + 1$. Mengapa jawabannya berbeda?
\end{obeactivity}

Listing~\ref{lst:ecc-gf2m} memuat keempat perkakas aritmetika polinomial:
penjumlahan sebagai XOR, perkalian dengan reduksi berjalan, pembagian panjang,
dan pemeriksa ketereduksian.

\lstinputlisting[
    language=Python,
    caption={Penjumlahan, perkalian, pembagian, dan pemeriksaan ketereduksian
    polinom di dalam $GF(2^m)$},
    label={lst:ecc-gf2m},
    firstline=47,
    lastline=87
]{\codefile{python/ecc_uji.py}}

\begin{obeactivity}
\textbf{Aktivitas 14.3: Enumerasi dan Pelelaran Titik pada Kurva $GF(p)$}
1. Jalankan fungsi enumerasi pada Listing~\ref{lst:ecc-kurva} untuk kurva
   $y^2 \equiv x^3 + x + 6 \pmod{11}$. Bandingkan daftar titik yang dihasilkan
   dengan Tabel~\ref{tab:enumerasi-gf11} dan sebaran titik pada
   Gambar~\ref{fig:sebaran-titik-gf11}.
2. Ulangi untuk $y^2 \equiv x^3 + x + 1 \pmod{23}$. Berapa banyak titik
   berhingga yang ditemukan, dan titik mana yang biasanya terlewat bila
   pencacahan dilakukan dengan tangan?
3. Hitung $P(2,4) + Q(5,9)$ dan $2P$ dengan fungsi penjumlahan dan penggandaan.
   Bandingkan hasilnya dengan Contoh~\ref{ex:gf11}.
4. Jalankan fungsi orde untuk titik $P(2,4)$ pada kurva mod 11 dan untuk
   $B(3,10)$ pada kurva mod 23. Bandingkan hasilnya dengan orde grup pada
   Tabel~\ref{tab:pelelaran-gf11} dan dengan catatan koreksi pada
   Subbagian~\ref{subsec:enumerasi-gf-p}.
5. Amati fungsi perkalian skalar: hitung $23P$ dan hitung berapa kali fungsi
   penjumlahan dan penggandaan dipanggil.
\end{obeactivity}

Listing~\ref{lst:ecc-kurva} memuat aritmetika kurva di atas medan berhingga:
pemeriksaan keanggotaan titik, pencarian lawan, penjumlahan, penggandaan,
perkalian skalar, enumerasi, dan pencarian orde titik.

\lstinputlisting[
    language=Python,
    caption={Penjumlahan, penggandaan, enumerasi, dan orde titik pada kurva
    di atas $GF(p)$},
    label={lst:ecc-kurva},
    firstline=91,
    lastline=172
]{\codefile{python/ecc_uji.py}}

\begin{obeactivity}
\textbf{Aktivitas 14.4: Penjumlahan Titik pada Kurva di Atas Bilangan Riil}
1. Gunakan kurva $y^2 = x^3 + 2x + 4$ dengan $a = 2$ dan $b = 4$.
2. Diberikan titik $P(2, 4)$ dan $Q(0, 2)$, hitunglah $R = P + Q$ dengan
   fungsi pada Listing~\ref{lst:ecc-real}, lalu hitung pula $2P$.
3. Bandingkan hasilnya dengan Contoh~\ref{ex:real}. Mengapa hasil
   penggandaannya menjadi pecahan, sedangkan hasil penjumlahannya tetap
   bilangan bulat?
4. Gambarkan sendiri langkah geometris penggandaan titik mengikuti
   Gambar~\ref{fig:penggandaan-titik}. Mengapa kurva seperti ini tidak dipakai
   untuk kriptografi?
\end{obeactivity}

Program pada Listing~\ref{lst:ecc-real} memakai pecahan eksak, bukan bilangan
pecahan mengambang. Dengan cara itu hasilnya tidak terpengaruh pembulatan.

\lstinputlisting[
    language=Python,
    caption={Penjumlahan dan penggandaan titik pada kurva di atas bilangan
    riil dengan pecahan eksak},
    label={lst:ecc-real},
    firstline=197,
    lastline=217
]{\codefile{python/ecc_uji.py}}

\begin{obeactivity}
\textbf{Aktivitas 14.5: Pengkodean Pesan dengan Metode Koblitz}
1. Jalankan fungsi pengkode pada Listing~\ref{lst:ecc-koblitz} untuk huruf
   \texttt{B} ($m = 11$) pada kurva $y^2 \equiv x^3 - x + 188 \pmod{751}$
   dengan parameter basis $k = 20$.
2. Periksa bahwa nilai $x$ yang dicoba berturut-turut adalah $221$, $222$,
   $223$, dan $224$, dan bahwa hanya yang terakhir memberi titik.
3. Tuliskan titik yang diperoleh dan kedua nilai ordinatnya. Periksa bahwa
   titik itu benar-benar memenuhi kurva.
4. Jalankan fungsi pendekode untuk memulihkan $m$ dari absis titik itu
   (Persamaan~\eqref{eq:koblitz-dekode}). Apakah hasilnya kembali menjadi
   $11$?
5. Ulangi seluruh langkah untuk huruf \texttt{K} dan \texttt{R}. Berapa
   percobaan yang diperlukan masing-masing?
\end{obeactivity}

Listing~\ref{lst:ecc-koblitz} memuat kedua sisi pemetaan: pencarian nilai $x$
yang menghasilkan titik, dan pemulihan nilai $m$ dari absis titik.

\lstinputlisting[
    language=Python,
    caption={Pengkodean pesan menjadi titik dan pendekodeannya dengan metode
    Koblitz},
    label={lst:ecc-koblitz},
    firstline=221,
    lastline=239
]{\codefile{python/ecc_uji.py}}

\begin{obeactivity}
\textbf{Aktivitas 14.6: Pertukaran Kunci ECDH dan Enkripsi ECEG}
1. Jalankan program pada Listing~\ref{lst:ecc-protokol} untuk kurva
   $y^2 \equiv x^3 + 2x + 1 \pmod{5}$ dengan titik basis $B(0,1)$, kunci
   privat Alice $a = 2$, dan kunci privat Bob $b = 3$. Periksa hasilnya
   terhadap Tabel~\ref{tab:ecdh-gf5}.
2. Periksa bahwa kedua arah menghitung kunci bersama yang sama, sesuai
   Persamaan~\eqref{eq:ecdh-kunci} dan Contoh~\ref{ex:ecdh}.
3. Jalankan bagian ECEG dengan kunci privat penerima $b = 3$, plainteks
   $P_M(7,12)$, dan kunci sesaat $k = 5$. Tuliskan cipherteksnya.
4. Pulihkan plainteksnya mengikuti Persamaan~\eqref{eq:eceg-dekripsi}. Apakah
   hasilnya kembali menjadi $P_M(7,12)$?
5. Ulangi enkripsi dengan kunci sesaat $k = 9$ dan bandingkan cipherteksnya
   dengan yang pertama. Plainteksnya sama, tetapi mengapa cipherteksnya
   berbeda?
\end{obeactivity}

Listing~\ref{lst:ecc-protokol} menjalankan kedua protokol. Perhatikan bahwa
kurva mod 5 sengaja dipilih kecil agar seluruh langkah dapat diperiksa dengan
tangan.

\lstinputlisting[
    language=Python,
    caption={Pertukaran kunci ECDH dan enkripsi--dekripsi ECEG},
    label={lst:ecc-protokol},
    firstline=364,
    lastline=411
]{\codefile{python/ecc_uji.py}}

Seluruh blok pengujian pada program ini dapat dijalankan sekaligus dengan
\texttt{python ecc\_uji.py}. Bila setiap angka pada bab ini cocok, program akan
mencetak \texttt{LULUS} pada setiap blok pengujian. Program kalkulator kurva
eliptik yang lebih lengkap tersedia di \texttt{code/python/ecc.py}.
